\chapter[El Cordón Cortado]{El Cordón Cortado\\ \large Quien olvida sus raíces marcha perdido en un mundo sin luz.}

\elegantimage{17_orfandad_filial/web/assets/hero.jpg}

\lettrine[lines=3, nindent=0.5em, findent=0.2em]{L}{}as raíces de nuestra encarnación son, por derecho divino, aquellos que nos han prestado su sangre. Quien voltea el rostro, reniega, o maldice a los padres que lo criaron, está talando el propio soporte sutil de su existencia...

La ingratitud filial corta indiscriminadamente aquel cordón dorado que nos ancla a la compasión cósmica. Es la afrenta capital contra aquellos que nos arroparon cuando no éramos más que llanto e indefensión.

Por lo tanto, la justicia universal es drástica con quien desmiembra la familia natural: arroja invariablemente al infractor al próximo nacimiento sumido en la orfandad absoluta, el abandono y el llanto silencioso de no tener jamás unos brazos que lo contengan.

\section{El Árbol sin Raíces}
\begin{elegantquote}
\textit{Ataviado con ropas de seda y jactándose de haber 'forjado su propio destino', un joven adinerado cerró la pesada puerta a sus padres ancianos, riendo ante sus ruegos por cobijo, y arguyendo que ya no le eran útiles...}

\textit{Lo que el joven soberbio ignoraba, es que el universo pesa las lágrimas de los padres sobre los hijos. Y al romper el vínculo por desdén, el telar cortó de tajo su continuidad protectora.}

\textit{Cuando exhaló su último aliento y volvió a nacer, su primer respiro fue bajo la tormenta fría, abandonado cruelmente entre las piedras. Así, descubrió en la intemperie desgarradora, el valor infinito que encierra poder pronunciar la palabra 'madre'.}

\end{elegantquote}

\elegantimage{17_orfandad_filial/web/assets/art.jpg}

\vspace{1em}
\begin{center}
    \textbf{\Large Rechazar a quien te dio la vida desgarra los hilos de tu propio destino.}
\end{center}
\vspace{1em}

\section{Análisis Kármico y Traducción}
\begin{wrapfigure}{r}{0.5\textwidth}
  \vspace{-1.5em}
  \begin{center}
  \inlineelegantimage[0.45\textwidth]{17_orfandad_filial/web/assets/pasaje_original.jpg}
  \end{center}
  \vspace{-1.5em}
\end{wrapfigure}
Quienes nos han prestado su sangre para encarnar en este pasaje existencial son el sagrado cordón ineludible que nos vincula a la raíz cósmica. La ingratitud filiar o el abandono frío a los que forjaron nuestra existencia rompe drásticamente esa protección vital natural. Ese desdén arroja invariablemente al individuo responsable hacia la brutalidad de la indefensión máxima y la orfandad en frío abandono para su próxima rueda experiencial.

